% !PRINTED 198
\weberpage{198}
%<<CONT
equal to $2n-2$. These values shall be called the \gesperrt{poles of the
group}. The exact enumeration of these poles gives us the
most important information about the number and nature of the
possible finite groups.

If the substitution coefficients $\beta, \gamma$ are equal to zero, then
$\Theta$ is a multiplicative substitution
\[
\begin{pmatrix} \alpha, & 0 \\ 0, & \delta \end{pmatrix}.
\]

In order that this substitution can be of finite degree,
the quotient $\alpha : \delta = \varepsilon$ must be a \gesperrt{root of unity
whose degree is a divisor of $n$}. As poles of this substitution
we regard $x = 0$ and $x = \infty$, both of which satisfy the following condition:
\[
x = \varepsilon x
\]

By (4) one can derive from $G$ an isomorphic group
\[
L G L^{-1} = G'
\]
if for $L$ an arbitrary linear substitution
\[
L = \begin{pmatrix} A, & B \\ C, & D \end{pmatrix}
\]
is taken. The poles of this group are obtained by applying
the substitution $L$ to the poles of $G$. If, then,
$a$ is one of the poles of $G$, then
\[
a' = \frac{Aa + B}{Ca + D}
\]
is the corresponding pole of $G'$.

One can determine the substitution coefficients $A, B, C, D$ in such a way
that three of the poles of $G'$ receive prescribed values.
In order, for example, to make the poles $0, \infty, 1$ of $G'$ correspond
to the poles $a, b, c$ of $G$, one sets
\[
L(x) = \frac{c-b}{c-a} \, \frac{x-a}{x-b}.
\]
If $a$ and $b$ are the poles of one and the same substitution $\Theta$
of $G$, then the corresponding substitution of $G'$ is accordingly
multiplicative.

From this we obtain the theorem:

\par\noindent 2.\ \gesperrt{For every group of linear substitutions one
can find a transformed group in which any one
of the given substitutions except the identity}
%CONT>>
